% Physics of Fluids (AIP Publishing), REVTeX 4.2 with the aip and pof options.
% First draft, 22 November 2026. The pof substyle is single column in REVTeX 4.2.
\documentclass[aip,pof,reprint,amsmath,amssymb]{revtex4-2}

\usepackage[utf8]{inputenc}
\usepackage[T1]{fontenc}
\usepackage{lmodern}
\usepackage{bm}
\usepackage{graphicx}
\usepackage{booktabs}
\usepackage{placeins}
\usepackage[colorlinks=true,linkcolor=blue,citecolor=blue,urlcolor=blue]{hyperref}
\usepackage{xurl}

\newcommand{\figfile}[2][width=\linewidth]{%
  \IfFileExists{#2}{\includegraphics[#1]{#2}}{%
    \fbox{\parbox[c][4cm][c]{0.8\linewidth}{\centering Missing figure:\\
      \texttt{\detokenize{#2}}}}}}
\newcommand{\bu}{\bm{u}}
\newcommand{\bE}{\bm{E}}
\newcommand{\dd}{\,\mathrm{d}}
\newcommand{\jump}[1]{[\![#1]\!]}

\begin{document}
\makeatletter\immediate\write\@auxout{\string\citation{aipcontrol}}\makeatother

\title{Oscillatory onset of electroconvection in two stacked layers of dielectric liquids}

\author{Pedro A. V\'azquez}
\email{pvazquez@us.es}
\affiliation{Department of Applied Physics III, ETSI, Universidad de Sevilla,
Camino de los Descubrimientos s/n, 41092 Sevilla, Spain}
% TODO: co-authors to be decided.

\date{First draft, 22 November 2026}

\begin{abstract}
Two layers of miscible dielectric liquids of different density, stacked between plane
electrodes, become unstable when charge is injected from the lower electrode. Numerical
simulations of this configuration show a threshold slightly above that of a single liquid
and ripples near the interface at the onset of convection. Starting from the equations of
two miscible liquids, we take the limit of infinite Schmidt number. The interface is then a
sharp material surface, its jump conditions follow from the conservation laws, and the
base state is steady. The interface conditions include gravity, the Maxwell stresses and
the conservation of current. A kinematic argument shows that, once the interface feeds back
on the flow, a stationary mode cannot carry fluid across it. With equal electrical
properties the onset of a single roll spanning both layers is therefore a Hopf
bifurcation for any stable stratification. The critical electric Rayleigh number grows
linearly with the Grashof number for weak stratification and the frequency grows as its
square root. Above a Grashof number of about 2140 a stationary mode made of two rows of
co-rotating cells, separated by a motionless interface, becomes critical. Removing fluid
inertia leaves the threshold and the frequency almost unchanged, while removing the
transit-time lag of the injected charge suppresses the dependence on stratification. The
ripples are thus transit-time oscillations of the charge, restored by gravity, and not an
inertia--gravity oscillation. The threshold reported in the simulations corresponds to a
Grashof number of about 13. The results are obtained with a two-domain Chebyshev method and
confirmed with a finite-element solver, both through its eigenvalues and through time
integration.
\end{abstract}

\maketitle

%===========================================================================
\section{Introduction}
%===========================================================================

A dielectric liquid between two electrodes becomes unstable when charge is injected from
one of them and the electric Rayleigh number $T$ exceeds a critical value. The linear
threshold of this electroconvection is classical. For strong injection it tends to
$T_c\approx 160$, and for an injection level $C=10$ it is $T_c=164.1$
\cite{AttenMoreau1972,Lacroix1975,Castellanos1998,Perez2014}. The instability is
subcritical, and the motion is sustained down to a lower nonlinear threshold.

Traor\'e \emph{et al.} \cite{Traore2024} studied a related configuration in which two
miscible dielectric liquids fill the gap, one above the other. The lower liquid is the
denser one, so the layer is stably stratified. Electroconvection mixes the two liquids
much faster than molecular diffusion. Their simulations show two features that the
single-liquid problem does not have. The forward threshold of convection is slightly
higher than that of a single liquid ($T=167.8$ instead of $164$ for $C=10$). In addition,
ripples appear near the initial interface during the growth of the instability. A
semi-analytical model based on a single prescribed roll attributed the ripples to a
competition between fluid inertia and gravity.

These observations raise three questions. Why does the threshold depend on the
stratification? What sets the frequency of the ripples? And what is the structure of the
mode that becomes unstable? We answer them with a linear stability analysis in the limit of
infinite Schmidt number. In this limit the concentration does not diffuse, the interface
is sharp and the base state is steady, so the problem has normal modes. This distinguishes
it from the simulations, where the interface spreads by diffusion and the base state
depends on time.

The interface displacement $\eta$ is governed by the kinematic condition
$\partial_t\eta=v$. A stationary mode therefore has no velocity across the interface. This
simple constraint organizes the whole problem. We show that, for equal electrical
properties, any stable stratification turns the onset of a single roll into a Hopf
bifurcation. The threshold increase is then set by the lag between the velocity field and
the injected charge, which needs a finite time to cross the layer. The role of this
transit time also governs the oscillatory modes of the single-liquid problem
(P. A. V\'azquez and A. T. P\'erez, in preparation). Here it becomes the mechanism that
controls the onset.

The paper is organized as follows. Section~\ref{sec:formulation} starts from the
dimensional equations of two miscible liquids, takes the limit of infinite Schmidt number,
derives the interface conditions from the conservation laws and linearizes them about the
base state. Section~\ref{sec:numerics} describes the
numerical methods and their validation, including a finite-element verification. The
results for equal electrical properties are given in Sec.~\ref{sec:results}, and
Sec.~\ref{sec:discussion} discusses unequal properties, finite Schmidt numbers and the
comparison with the simulations of Ref.~\onlinecite{Traore2024}.

%===========================================================================
\section{Formulation}\label{sec:formulation}
%===========================================================================

\subsection{Dimensional equations}

Two miscible dielectric liquids fill the gap $0<y^*<H$ between plane electrodes
(Fig.~\ref{fig:sketch}). The lower electrode, at potential $V_0$, injects charge of
density $q_0$. The upper electrode is grounded. The problem is periodic in $x^*$, and
asterisks denote dimensional variables. The local composition is described by the volume
fraction $c$ of liquid~1, which is the reference liquid. Initially liquid~1 fills the upper
half of the gap, $c=1$ for $y^*>H/2$, and the denser liquid~2 fills the lower half, $c=0$.

\begin{figure}
\figfile[width=0.55\linewidth]{figures/fig_sketch.pdf}
\caption{Configuration. The lower liquid is the denser one. Charge is injected from the
lower electrode, and the interface is displaced by $\eta(x,t)$.}
\label{fig:sketch}
\end{figure}

The density, permittivity, ionic mobility and viscosity depend on the composition. We write
them as $\rho(c)$, $\varepsilon(c)$, $K(c)$ and $\mu(c)$, with the values of liquid~1 at
$c=1$ and those of liquid~2 at $c=0$. A linear mixing law between these values is the
simplest choice, but in the limit studied below only the values of the pure liquids
matter. In the Boussinesq approximation, with reference density $\rho_1$, the equations are
\begin{align}
&\nabla^*\!\cdot\bu^*=0,\label{eq:d_mass}\\
&\rho_1\left(\partial_{t^*}\bu^*+\bu^*\cdot\nabla^*\bu^*\right)=-\nabla^*p^*
+\nabla^*\!\cdot\left[\mu\left(\nabla^*\bu^*+\nabla^*\bu^{*T}\right)\right]\nonumber\\
&\qquad\qquad+q^*\bE^*-\tfrac12E^{*2}\nabla^*\varepsilon-\rho(c)\,g\,\hat{\bm y},
\label{eq:d_ns}\\
&\partial_{t^*}q^*+\nabla^*\!\cdot\left[q^*\left(\bu^*+K\bE^*\right)\right]=0,
\label{eq:d_charge}\\
&\nabla^*\!\cdot\left(\varepsilon\bE^*\right)=q^*,\qquad \bE^*=-\nabla^*V^*,
\label{eq:d_gauss}\\
&\partial_{t^*}c+\bu^*\cdot\nabla^*c=\mathcal{D}\,\nabla^{*2}c .\label{eq:d_conc}
\end{align}
The electric body force is the Korteweg--Helmholtz force without electrostriction, whose
contribution is a gradient absorbed in the pressure. With Eq.~\eqref{eq:d_gauss} it is the
divergence of the Maxwell stress,
\begin{equation}
q^*\bE^*-\tfrac12E^{*2}\nabla^*\varepsilon=\nabla^*\!\cdot\mathsf T^*_e,\qquad
\mathsf T^*_e=\varepsilon\,\bE^*\bE^*-\tfrac12\varepsilon E^{*2}\,\mathsf I .
\label{eq:d_maxwell}
\end{equation}
Charge diffusion is neglected, as usual for injection problems, and $\mathcal D$ is the
molecular diffusivity of the two liquids into each other. The boundary conditions are
$\bu^*=0$ on both electrodes, $V^*=V_0$ and $q^*=q_0$ at $y^*=0$, $V^*=0$ at $y^*=H$, and no
flux of $c$ through the walls.

\subsection{Dimensionless equations}

We scale lengths with $H$, time with the drift time $H^2/(K_1V_0)$, velocities with
$K_1V_0/H$, the potential with $V_0$, the charge density with $q_0$, the pressure with
$\rho_1K_1^2V_0^2/H^2$, and the properties with those of liquid~1. The dimensionless
properties are then $\varepsilon=K=\mu=1$ in liquid~1 and $\varepsilon_r$, $K_r$, $\mu_r$ in
liquid~2. The parameters of the classical problem are built with liquid~1,
\begin{equation}
T=\frac{\varepsilon_1V_0}{\mu_1K_1},\qquad C=\frac{q_0H^2}{\varepsilon_1V_0},\qquad
M=\frac{1}{K_1}\left(\frac{\varepsilon_1}{\rho_1}\right)^{1/2},
\end{equation}
as in Ref.~\onlinecite{Traore2024}. They give the ratios $\nu_1/(K_1V_0)=M^2/T$ of the
drift time $H^2/(K_1V_0)$ to the viscous time $H^2/\nu_1$ and $q_0H^2/(\rho_1K_1^2V_0)=CM^2$ of the Coulomb force to
inertia.

The buoyancy term needs one more step. With $\rho(c)=\rho_1+\Delta\rho\,(1-c)$ and
$\Delta\rho=\rho_2-\rho_1>0$, the gravity force per unit mass is
$-g\hat{\bm y}-(g\Delta\rho/\rho_1)(1-c)\hat{\bm y}$. Its uniform part is absorbed in the
pressure, and what remains is the buoyancy $(g\Delta\rho/\rho_1)\,c\,\hat{\bm y}$. In drift
units its coefficient is
\begin{equation}
G=\frac{g\Delta\rho\,H^3}{\rho_1K_1^2V_0^2}=Gr\,\frac{M^4}{T^2},\qquad
Gr=\frac{g\Delta\rho\,H^3}{\rho_1\nu_1^2}=\frac{Ra_S}{Sc},\qquad Sc=\frac{\nu_1}{\mathcal D},
\end{equation}
where $Gr$ is the Grashof number. It is the buoyancy parameter that appears when the
velocity is scaled with $\nu_1/H$, as in Ref.~\onlinecite{Traore2024}, while $G$ is its
counterpart in the drift scaling used here. The
dimensionless equations are
\begin{align}
\partial_t\bu+\bu\cdot\nabla\bu&=-\nabla p+\frac{M^2}{T}\nabla\cdot\left[\mu\left(\nabla\bu+
\nabla\bu^T\right)\right]+M^2\nabla\cdot\mathsf T_e+G\,c\,\hat{\bm y},\label{eq:ns}\\
\partial_t q+\nabla\cdot\left[q\left(\bu+K\bE\right)\right]&=0,\label{eq:charge}\\
\nabla\cdot\left(\varepsilon\bE\right)=Cq,\qquad \bE&=-\nabla\phi,\label{eq:poisson}\\
\partial_tc+\bu\cdot\nabla c&=\kappa\nabla^2c,\qquad \kappa=\frac{M^2}{T\,Sc},
\label{eq:conc}
\end{align}
with $\nabla\cdot\bu=0$, $\mathsf T_e=\varepsilon\bE\bE-\tfrac12\varepsilon E^2\mathsf I$ and
$M^2\nabla\cdot\mathsf T_e=CM^2q\bE-\tfrac12M^2E^2\nabla\varepsilon$. The boundary conditions
are $\bu=0$ on both electrodes, $\phi=1$ and $q=1$ at $y=0$, $\phi=0$ at $y=1$, and
$\partial_yc=0$ on both walls. Unless stated otherwise we take $C=M=10$.

\subsection{The limit of infinite Schmidt number}

In the simulations of Ref.~\onlinecite{Traore2024} the Schmidt number is large, and in real
liquids it is larger still. We consider the limit $Sc\to\infty$ at fixed $Gr$, so that
$\kappa=0$ while the buoyancy stays finite. Equation~\eqref{eq:conc} then states that $c$
is constant along fluid trajectories. Starting from the step profile, $c$ only takes the
values $0$ and $1$, and the two liquids remain separated by a material surface $S$,
$y=\tfrac12+\eta(x,t)$. The properties are constant on each side. In each layer the
equations reduce to those of a single liquid,
\begin{align}
\partial_t\bu+\bu\cdot\nabla\bu&=-\nabla p+\nu\nabla^2\bu+CM^2q\bE+G\,c\,\hat{\bm y},
\label{eq:ns_layer}\\
\partial_tq+\bu\cdot\nabla q+K\bE\cdot\nabla q&=-\frac{KC}{\varepsilon}\,q^2,
\label{eq:charge_layer}\\
\varepsilon\nabla^2\phi&=-Cq,\label{eq:poisson_layer}
\end{align}
with $\nu=\mu M^2/T$ and $c$ equal to $0$ or $1$.

The conditions on $S$ follow from integrating the conservation laws across a thin layer
around $S$ and letting its thickness go to zero. All fields stay bounded in that layer.
In particular the charge density is bounded, so the layer stores no surface charge. The
liquids are miscible, so there is no surface tension. Let $\bm n$ be the normal pointing
into the upper liquid, $V_n$ the normal velocity of $S$ and $\jump f$ the jump of $f$ from
the lower to the upper side. Then
\begin{align}
&\jump{\bu}=0, \qquad V_n=\bu\cdot\bm n, &&\text{(velocity, kinematic)}\label{eq:j_kin}\\
&\jump{\Sigma\cdot\bm n}=0, &&\text{(stress)}\label{eq:j_stress}\\
&\jump{\phi}=0,\qquad \jump{\varepsilon\bE\cdot\bm n}=0, &&\text{(potential, displacement)}
\label{eq:j_elec}\\
&\jump{Kq\,\bE\cdot\bm n}=0, &&\text{(current)}\label{eq:j_current}
\end{align}
with the total stress $\Sigma=-p\,\mathsf I+\nu\left(\nabla\bu+\nabla\bu^T\right)+M^2\mathsf
T_e$. The kinematic condition is Eq.~\eqref{eq:conc} with $\kappa=0$ applied to the
discontinuity of $c$. The velocity is continuous because the viscous stress is bounded.
The buoyancy is a bounded body force, so it does not enter Eq.~\eqref{eq:j_stress} directly.
The stress condition contains the Maxwell stress, and with it the polarization force
$-\tfrac12M^2E^2\nabla\varepsilon$, which is concentrated on $S$ when $\varepsilon$ jumps. In
the current condition the convective flux $q(\bu\cdot\bm n-V_n)$ vanishes on both sides by
the kinematic condition. These conditions are exact. They are linearized below.

\subsection{Base state}

The base state is at rest, with the interface flat. In each layer
$(\varepsilon\bar E)'=C\bar q$, and the current $J=K\bar q\bar E$ is uniform. With
$E_0=\bar E(0)$ and $J=K_rE_0$,
\begin{align}
\bar E^2&=E_0^2+\frac{2CE_0}{\varepsilon_r}\,y,&&0<y<\tfrac12,\\
\bar E^2&=E_i^2+2CK_rE_0\left(y-\tfrac12\right),&&\tfrac12<y<1,
\end{align}
with $E_i=\varepsilon_r\bar E(\tfrac12^-)$ and $\bar q=J/(K\bar E)$. The value of $E_0$
follows from $\int_0^1\bar E\dd y=1$. This is the analytical solution of
Ref.~\onlinecite{Traore2024} written in a different notation. By
Eqs.~\eqref{eq:j_elec}--\eqref{eq:j_current}, the normal displacement $\varepsilon\bar E$
and the current are continuous at the interface, while $\bar E$ jumps when
$\varepsilon_r\ne1$ and $\bar q$ jumps by the factor $K_r/\varepsilon_r$
(Fig.~\ref{fig:base}).

\begin{figure*}
\figfile[width=\textwidth]{figures/fig_base.pdf}
\caption{Base state for $C=10$. (a) Electric field. (b) Charge density. The interface is
at $y=1/2$.}
\label{fig:base}
\end{figure*}

In each layer the divergence of the Maxwell stress equals the Coulomb force,
$M^2(\tfrac12\varepsilon\bar E^2)'=M^2\varepsilon\bar E\bar E'=CM^2\bar q\bar E$. The
hydrostatic balance $\bar p'=CM^2\bar q\bar E+G\bar c$ then shows that the base normal stress
$\bar\Sigma_{yy}=-\bar p+\tfrac12M^2\varepsilon\bar E^2$ satisfies
\begin{equation}
\bar\Sigma_{yy}'=-G\,\bar c .
\label{eq:basestress}
\end{equation}
This identity is how gravity enters the normal-stress condition at the displaced
interface.

\subsection{Linear perturbations}

We perturb the base state with normal modes proportional to $e^{ikx+\lambda t}$, and keep
the notation $v$, $q$, $\phi$ for the amplitudes of the perturbations of the vertical
velocity, the charge density and the potential.

\paragraph*{Concentration.}
The linearized form of Eq.~\eqref{eq:conc} is
\begin{equation}
\lambda c'=-v\,\bar c'+\kappa\,(D^2-k^2)\,c',
\label{eq:lin_conc}
\end{equation}
with $D=\mathrm d/\mathrm dy$. For $\kappa=0$ and the step profile $\bar c=\Theta(y-\tfrac12)$, with $\Theta$ the Heaviside function,
$\bar c'=\delta(y-\tfrac12)$, and the solution is
\begin{equation}
c'=-\eta\,\delta\!\left(y-\tfrac12\right),\qquad \lambda\eta=v(\tfrac12).
\label{eq:kinematic}
\end{equation}
The concentration perturbation is a displacement of the interface. Its only degree of
freedom is the amplitude $\eta$, and its evolution equation is the linearized kinematic
condition. The buoyancy $G\,c'\,\hat{\bm y}$ is then a force concentrated on the interface.
It appears in the normal-stress condition below, through Eq.~\eqref{eq:basestress}, and not
in the bulk equations.

\paragraph*{Bulk equations.}
In each layer, with $\zeta=(D^2-k^2)v$, Eqs.~\eqref{eq:ns_layer}--\eqref{eq:poisson_layer}
give
\begin{align}
\lambda\zeta&=\nu(D^2-k^2)\zeta-k^2CM^2\left(\bar E q+\bar q'\phi\right),\label{eq:lin_mom}\\
\lambda q&=-\bar q'v-K\bar E\,Dq-\frac{2KC}{\varepsilon}\,\bar q\,q+K\bar q'\,D\phi,
\label{eq:lin_charge}\\
0&=\varepsilon(D^2-k^2)\phi+Cq .\label{eq:lin_poisson}
\end{align}
Equation~\eqref{eq:lin_mom} is the curl of the momentum equation, with the pressure
eliminated. At the walls $v=Dv=0$ and $\phi=0$, and $q=0$ at the injector. The charge
equation is hyperbolic, so no condition is imposed on $q$ at the collector.

\paragraph*{Interface conditions.}
The conditions \eqref{eq:j_kin}--\eqref{eq:j_current} hold on the displaced interface.
We transfer them to $y=\tfrac12$ with $f(\tfrac12+\eta)\simeq\bar f(\tfrac12)+\hat f(\tfrac12)+
\eta\bar f'(\tfrac12)$, where $\hat f$ is the perturbation, with normal $\bm n\simeq(-ik\eta,1)$ and tangent
$\bm t\simeq(1,ik\eta)$. Velocity continuity gives $\jump v=0$ and, with the continuity
equation $iku+Dv=0$, $\jump{Dv}=0$. The kinematic condition is
Eq.~\eqref{eq:kinematic}. Continuity of the potential, $\jump{\bar\phi(\tfrac12+\eta)+\phi}=0$,
and of the normal displacement give
\begin{equation}
\jump\phi=\eta\,\jump{\bar E},\qquad \jump{\varepsilon D\phi}=\eta\,\jump{(\varepsilon\bar E)'}
=\eta\,C\jump{\bar q}.
\label{eq:electric}
\end{equation}
The current condition gives $\jump{K(\bar qE_y+q\bar E)}+\eta\jump{(K\bar q\bar E)'}=0$, and
the second term vanishes because the base current is uniform. Hence
\begin{equation}
\jump{K(\bar Eq-\bar q\,D\phi)}=0 .
\label{eq:current}
\end{equation}
The tangential stress $\bm t\cdot\Sigma\cdot\bm n$ has a viscous part,
$\jump{\nu(Du+ikv)}$, and an electric part,
$M^2\jump{\varepsilon\bar E\,(E_x+ik\eta\bar E)}=ikM^2\varepsilon\bar E\,(-\jump\phi+
\eta\jump{\bar E})$. The electric part vanishes by Eq.~\eqref{eq:electric}, because
$\varepsilon E_n$ and $E_t$ are both continuous. The shear condition is then
\begin{equation}
\jump{\nu(\zeta+2k^2v)}=0 .
\end{equation}
The normal stress gives, with Eq.~\eqref{eq:basestress} and the pressure taken from the
horizontal momentum equation (Appendix~\ref{app:normal}),
\begin{equation}
\jump{-\nu D\zeta+2\nu k^2Dv}+k^2CM^2\jump{\bar q\phi}-k^2M^2\jump{\varepsilon\bar E\,D\phi}
=k^2G\,\eta .
\label{eq:normalstress}
\end{equation}
The problem has fourteen constants of integration in the two layers and the unknown $\eta$.
They are fixed by seven wall conditions and eight interface conditions, namely
$\jump v=\jump{Dv}=0$, the shear and normal stresses, the kinematic condition, the two
electric conditions \eqref{eq:electric} and the current condition \eqref{eq:current}.

\subsection{A kinematic constraint}\label{sec:kinematic}

For a stationary mode ($\lambda=0$), Eq.~\eqref{eq:kinematic} requires $v(\tfrac12)=0$.
The interface can be displaced, but no fluid crosses it. This constraint matters whenever
$\eta$ acts back on the flow. That happens through the normal stress when $Gr\ne0$, through
the potential jump when $\varepsilon_r\ne1$ and through the displacement condition when
$K_r\ne\varepsilon_r$. The viscosity ratio never involves $\eta$. Only for $Gr=0$ and
$\varepsilon_r=K_r=1$ is the interface a passive marker.

Two consequences follow. A single roll spanning both layers can only become unstable
through a Hopf bifurcation. A stationary onset requires two rows of cells, one in each
layer. The situation is analogous to double-diffusive convection when the stabilizing
component does not diffuse \cite{Turner1973}.

The constraint is a property of the limit $Sc\to\infty$. For finite $Sc$ a stationary mode
only requires $v\,\bar c'=\kappa(D^2-k^2)c'$ in Eq.~\eqref{eq:lin_conc}. A vertical velocity
across the interface can then be balanced by diffusion of the concentration.

%===========================================================================
\section{Numerical methods}\label{sec:numerics}
%===========================================================================

\subsection{Two-domain Chebyshev collocation}

Each layer is discretized by Chebyshev collocation \cite{Trefethen2000} with $N+1$ points,
and the two layers are coupled through the interface conditions. The generalized
eigenvalue problem $A\bm x=\lambda B\bm x$ is solved with the QZ algorithm. Two choices are
essential for accuracy.

The first concerns the coordinate. With strong injection $E_0$ is small ($E_0=0.111$ for
$C=10$). The field $\bar E(y)$ then has a branch point at $y^*=-E_0/2C=-0.0055$, just
below the injector. The coefficient $K\bar E$ of the charge equation vanishes there, and
the eigenfunctions inherit the singularity. We collocate in the local field $s=\bar E(y)$,
that is $y=y_0+(s^2-a)/b$ in a layer where $\bar E^2=a+b(y-y_0)$. In this coordinate
$K\bar E\,\partial_y=(Kb/2)\,\partial_s$. The charge is advected at constant speed, and
the coefficients are regular.

The second concerns the momentum equation. We keep $\zeta$ as an unknown instead of
collocating $D^4$, whose condition number grows like $N^8$. At each wall the
$\zeta$-definition row carries $v=0$ and the momentum row carries $Dv=0$. At the interface
the same rows of each layer carry the four velocity and stress conditions.

With collocation in $y$ and $D^4$, the leading eigenvalue changes at the $10^{-3}$ level
when the matrix entries are perturbed at the $10^{-14}$ level. Chebyshev collocation of an
advection operator with an inflow condition is exponentially non-normal
\cite{TrefethenTrummer1987}, and the fourth-order operator amplifies the effect. With both
choices the same test changes the eigenvalue by about $10^{-8}$, and no spurious mode
appears. Table~\ref{tab:convergence} shows the convergence of the leading eigenvalue. We
use $N=24$ per layer and keep only eigenvalues that persist at $N=32$.

\begin{table}
\caption{Leading eigenvalue for $C=M=10$, $T=159.73$, $k=2\pi/1.228$, equal properties and
$Gr=0$, versus the number of points per layer.}
\label{tab:convergence}
\begin{ruledtabular}
\begin{tabular}{rc}
$N$ & $\sigma$\\\hline
12 & $-0.1231781$\\
16 & $-0.1235440$\\
20 & $-0.1235831$\\
24 & $-0.1235870$\\
32 & $-0.1235874$\\
60 & $-0.1235874$\\
\end{tabular}
\end{ruledtabular}
\end{table}

\subsection{Validation}

Table~\ref{tab:validation} summarizes the tests. In the homogeneous limit the method
reproduces the classical thresholds. Without space charge ($C\to0$) the interface supports
damped gravity waves. Their frequency tends to the inviscid value
$\omega^2=Gk/[\coth(k/2)+\coth(k/2)]$ with a correction of order $\nu^{1/2}$, the Stokes
layer. For $Gr<0$ the Rayleigh--Taylor instability appears. With $\varepsilon_r\ne1$ and no
space charge the interface is destabilized by the normal field. For equal depths the
neutral condition is
\begin{equation}
G=M^2k\coth(k/2)\,\frac{\varepsilon_r(1-\varepsilon_r)^2}{1+\varepsilon_r}\,E_L^2,\qquad
E_L=\frac{2}{1+\varepsilon_r},
\end{equation}
in agreement with Melcher \cite{Melcher1981}. These tests check the kinematic, gravity,
potential-jump, displacement-jump and Maxwell normal-stress terms. The terms specific to
$K_r\ne\varepsilon_r$ are only checked through the smooth-profile solver described below.

\begin{table}
\caption{Validation of the two-domain solver.}
\label{tab:validation}
\begin{ruledtabular}
\begin{tabular}{lll}
Test & Result & Reference\\\hline
$T_c$, $C=5$ & 175.7101, $k_c=5.0413$ & 175.7, 5.041\\
$T_c$, $C=10$ & 164.0614, $k_c=5.1125$ & 164.1 \cite{Perez2014}\\
Gravity waves, $\nu=10^{-3},10^{-4},10^{-5}$ & $\omega/\omega_{\rm inv}=0.918,\,0.974,\,0.992$ & $1-O(\nu^{1/2})$\\
Normal field, $\varepsilon_r=2$ & neutral $G$ within $2\times10^{-7}$ & Ref.~\onlinecite{Melcher1981}\\
\end{tabular}
\end{ruledtabular}
\end{table}

\subsection{Smooth interface}

Simulations cannot represent a sharp interface. Without diffusion, however, any
concentration profile $\bar c(y)$ at rest is an exact steady state. We therefore also
solve the problem with
\begin{equation}
\bar c=\tfrac12\left[1+\tanh\left(\frac{y-\tfrac12}{\delta}\right)\right]
\end{equation}
and equal properties. For $\kappa=0$ the linearized concentration equation
\eqref{eq:lin_conc} gives $c'=-\bar c'\xi$, with $\lambda\xi=v$ the vertical displacement,
and the buoyancy adds the term $k^2G\,\bar c'\xi$ to Eq.~\eqref{eq:lin_mom}. The two domains are joined by plain continuity conditions. As
$\delta\to0$, $\bar c'\to\delta(y-\tfrac12)$ and $\xi(\tfrac12)\to\eta$. The neutral point
converges to the sharp-interface value as $\delta^2$ (Fig.~\ref{fig:delta}), which is an
independent check of the interface conditions.

Without diffusion the stratified fluid also has slow buoyancy modes. Their eigenvalues
accumulate at $\lambda=0$, with exact zeros where $\bar c'$ underflows. The largest real
part is therefore not a valid selector of the electroconvective mode, and we follow the
oscillatory mode explicitly.

\begin{figure}
\figfile[width=0.55\linewidth]{figures/fig_delta.pdf}
\caption{Relative difference between the neutral electric Rayleigh number of the smooth
profile of thickness $\delta$ and that of the sharp interface, for $k=5$.}
\label{fig:delta}
\end{figure}

\subsection{Finite-element verification}

The finite-element solver of our study of the single-liquid problem uses
Firedrake \cite{Rathgeber2016}. It is written in perturbation variables, with the
residual of the perturbation defined as the difference between the residual of the total
fields and that of the base state. The zero perturbation is then an exact discrete
solution. We add the concentration field, governed by Eq.~\eqref{eq:conc} with
$\kappa=0$, to the mixed space (Q2/Q1/Q1/Q2/Q2 for $\bu,p,q,\phi,c$) and the buoyancy force
to the momentum equation, with $\delta=0.02$ and
a mesh clustered near the interface. The domain is periodic with one wavelength
($k=5$).

The temporal eigenvalues are those of $-J\bm x=\lambda M\bm x$, where $J$ and $M$ are the
Jacobian and the time-derivative matrix of the residual of the time stepper. Because of
the cluster of buoyancy modes at $\lambda=0$, we use shift-invert Arnoldi with a complex
shift close to the expected $i\omega$. The neutral $T$ follows from a secant iteration on
$\mathrm{Re}\,\lambda$.

Table~\ref{tab:firedrake} compares the results on a mesh of $62\times100$ elements with the
Chebyshev values for the same profile. The frequencies agree within 0.3\%. The thresholds
differ by $-0.29\%$ in the weakly stratified case. This is the discretization offset of
the mesh, which is coarser in $x$ than our production mesh for the single liquid. Measured
from that reference, the shift of the threshold caused by the stratification agrees within
2\%.

\begin{table}
\caption{Neutral electric Rayleigh number and frequency (scaled with the drift time) of the Hopf mode for
$\delta=0.02$ and $k=5$, from the finite-element eigenproblem ($62\times100$) and from the
Chebyshev solver. $\Delta T_n=T_n(Gr)-T_n(0.1)$.}
\label{tab:firedrake}
\begin{ruledtabular}
\begin{tabular}{rcccccc}
 & \multicolumn{3}{c}{Finite elements} & \multicolumn{3}{c}{Chebyshev}\\
$Gr$ & $T_n$ & $\omega$ & $\Delta T_n$ & $T_n$ & $\omega$ & $\Delta T_n$\\\hline
0.1 & 163.735 & 0.1005 & --- & 164.204 & 0.1006 & ---\\
10 & 166.583 & 0.9863 & 2.848 & 167.002 & 0.9872 & 2.798\\
100 & 191.553 & 2.6826 & 27.818 & 191.677 & 2.6874 & 27.473\\
1000 & 356.213 & 4.2387 & 192.48 & 355.631 & 4.2494 & 191.43\\
\end{tabular}
\end{ruledtabular}
\end{table}

The eigenvalues are confirmed by time integration of the full nonlinear perturbation
equations (BDF2, $\Delta t=0.02$) from a small charge seed. The growth rate and frequency
are fitted with the matrix-pencil method \cite{HuaSarkar1990} to the projection $a_c(t)$
of the vertical velocity on the box wavenumber (Fig.~\ref{fig:firedrake}). On a mesh of
$32\times40$ elements at $Gr=100$ and $T=192$, the eigenproblem gives
$\lambda=0.024379+2.642578i$. The time integration gives $0.024217\pm2.640130i$ for a
seed of norm $E_0=10^{-8}$ and $0.024214\pm2.640131i$ for $E_0=10^{-12}$. The independence
of the amplitude confirms the linear regime. The remaining difference is the time error of
BDF2 at $\omega\Delta t=0.053$.

\begin{figure}
\figfile[width=0.55\linewidth]{figures/fig_firedrake.pdf}
\caption{Finite-element time integration at $Gr=100$, $T=192$, on a $32\times40$ mesh.
Projection of the vertical velocity on the box wavenumber for two seed amplitudes, with the
envelope given by the fitted growth rate.}
\label{fig:firedrake}
\end{figure}

%===========================================================================
\section{Results for equal electrical properties}\label{sec:results}
%===========================================================================

\subsection{Neutral curves and thresholds}

Figure~\ref{fig:neutral} shows the neutral curves for several Grashof numbers. For
$Gr>0$ the leading mode at small and moderate $k$ is oscillatory. The neutral curve
rises with $Gr$ and its minimum moves to smaller $k$. At large $k$ the leading mode is
stationary and its neutral curve does not depend on $Gr$. This is the two-tier mode of
Sec.~\ref{sec:kinematic}. With equal properties and $v(\tfrac12)=0$, $\eta$ is decoupled
from the electric problem, so gravity cannot act on this mode.

\begin{figure}
\figfile[width=0.55\linewidth]{figures/fig_neutral.pdf}
\caption{Neutral curves of the leading mode. For $Gr>0$ only the oscillatory part is
drawn. The dashed curve is the stationary two-tier mode, which does not depend on $Gr$.}
\label{fig:neutral}
\end{figure}

The critical values are listed in Table~\ref{tab:critical} and shown in
Fig.~\ref{fig:TcGr}. For any $Gr>0$ the onset is a Hopf bifurcation. For weak
stratification
\begin{equation}
T_c-164.06\simeq0.295\,Gr,\qquad \omega_c\simeq0.317\,Gr^{1/2},
\end{equation}
with $\omega_c$ scaled with the drift time. The stationary two-tier mode has $T_c=477.68$ at
$k_c=10.42$, about twice the wavenumber of the full-depth roll, as expected for cells of
half the height. It becomes critical for $Gr>Gr^*=2137.6$. At $Gr^*$ the two modes have
very different wavenumbers ($k=3.60$ for the oscillatory one).

\begin{table}
\caption{Critical points for equal properties, $C=M=10$. The frequency is scaled with the
drift time. The last row is the stationary
two-tier mode, which does not depend on $Gr$.}
\label{tab:critical}
\begin{ruledtabular}
\begin{tabular}{rccc}
$Gr$ & $T_c$ & $k_c$ & $\omega_c$\\\hline
0 & 164.0614 & 5.1125 & 0\\
0.1 & 164.0909 & 5.1122 & 0.1005\\
1 & 164.3563 & 5.1093 & 0.3172\\
10 & 167.0006 & 5.0825 & 0.9863\\
100 & 192.0616 & 4.8743 & 2.6968\\
1000 & 352.5673 & 4.0438 & 4.5180\\
2137.6 & 477.6823 & 3.5954 & 4.7558\\\hline
two-tier & 477.6823 & 10.4165 & 0\\
\end{tabular}
\end{ruledtabular}
\end{table}

\begin{figure*}
\figfile[width=\textwidth]{figures/fig_Tc_Gr.pdf}
\caption{(a) Critical electric Rayleigh number versus the Grashof number for the full
model, without fluid inertia and without the transit-time lag of the charge. (b) Critical
frequency scaled with the viscous time $H^2/\nu_1$ for the same three models.}
\label{fig:TcGr}
\end{figure*}

\subsection{Structure of the modes}

On the whole oscillatory branch the critical mode is a single roll spanning both layers
(Fig.~\ref{fig:eigen}). The vertical velocity does not change sign across the depth and is
largest near the interface. The two-tier mode is made of two rows of cells rotating in the
same sense. At the interface both $v$ and $Dv$ vanish, so both velocity components are
zero there. The interface acts as an internal no-slip wall.

\begin{figure*}
\figfile[width=\textwidth]{figures/fig_eigen.pdf}
\caption{Critical eigenfunctions. (a) Modulus of the vertical velocity. (b) Streamlines of
the Hopf mode at $Gr=100$ at one phase. (c) Streamlines of the stationary two-tier mode.
Solid and dashed lines have opposite signs.}
\label{fig:eigen}
\end{figure*}

\subsection{Mechanism}

To identify the mechanism we remove, one at a time, the two sources of delay in the
problem. Without fluid inertia ($\lambda\zeta$ dropped) the velocity is slaved to the
forces. Without charge lag ($\lambda q$ dropped) the charge responds instantaneously to
the velocity. The interface keeps its own dynamics, Eq.~\eqref{eq:kinematic}, in both
cases. The results are given in Table~\ref{tab:mechanism} and Fig.~\ref{fig:TcGr}.

\begin{table}
\caption{Critical electric Rayleigh number and frequency (scaled with the drift time) for the full
model, without fluid inertia and without charge lag.}
\label{tab:mechanism}
\begin{ruledtabular}
\begin{tabular}{rcccccc}
 & \multicolumn{2}{c}{Full} & \multicolumn{2}{c}{No inertia} & \multicolumn{2}{c}{No charge lag}\\
$Gr$ & $T_c$ & $\omega_c$ & $T_c$ & $\omega_c$ & $T_c$ & $\omega_c$\\\hline
10 & 167.00 & 0.99 & 167.45 & 1.08 & 164.15 & 2.47\\
100 & 192.06 & 2.70 & 195.94 & 2.90 & 164.91 & 7.78\\
1000 & 352.57 & 4.52 & 363.36 & 4.58 & 172.37 & 23.41\\
\end{tabular}
\end{ruledtabular}
\end{table}

Without inertia both the threshold and the frequency are almost unchanged. Without charge
lag the threshold hardly depends on $Gr$, and the frequency becomes that of a much faster
inertia--gravity oscillation. The onset is therefore controlled by three ingredients.
The interface integrates the vertical velocity, gravity restores it, and the charge
responds to the flow with the delay of its transit time. For $M=10$ the viscous
relaxation of the flow is fast compared with this transit time, so inertia plays a
minor role.

A simple argument shows why the lag raises the threshold. In a one-roll model with
instantaneous charge, as in Ref.~\onlinecite{Traore2024}, the amplitude obeys a damped
oscillator equation whose damping changes sign at a value of $T$ independent of $Gr$.
Gravity only sets the frequency. When the charge at height $y$ is created by the velocity
at earlier times, the Coulomb coupling becomes a function of the growth rate,
\begin{equation}
a_2(\lambda)=\iint u_{0y}(y)\,u_{0y}(y')\,e^{-\lambda(y-y')\tau}\dd y'\dd y,
\end{equation}
with $\tau$ the transit time. At an oscillatory onset, $\lambda=i\omega$, the destabilizing
part is $\mathrm{Re}\,a_2(i\omega)<a_2(0)$. The threshold then grows with the frequency,
and hence with $Gr$.

%===========================================================================
\section{Discussion}\label{sec:discussion}
%===========================================================================

\subsection{Comparison with the simulations}

The forward threshold $T=167.8$ reported in Ref.~\onlinecite{Traore2024} for $C=10$
corresponds to $Gr=12.7$ in the present analysis. This value lies inside the range of
$Ra_S/Sc$ explored in those simulations. The critical frequency is then $1.107$ when time
is scaled with the drift time. With the viscous time $H^2/\nu_1$ the frequency is
$\omega_cT_c/M^2=1.86$ and the period is $3.4$. The stratification alone explains the
increase of the threshold, and the ripples are the Hopf mode of Sec.~\ref{sec:results}.

Two differences must be kept in mind. The simulations have a finite Schmidt number, and
the interface spreads by diffusion over a thickness of order $2(t/Sc)^{1/2}$. For
$Sc=1000$ and $t\approx100$ this is comparable to the layer thickness at the onset of
convection. Diffusion also allows stationary modes that carry fluid across the interface,
because a vertical velocity can be balanced by diffusion of the concentration. Second, the
backward branch of the bifurcation diagram of Ref.~\onlinecite{Traore2024} is computed
after the liquids have mixed, so it corresponds to a homogeneous liquid.

\subsection{Unequal electrical properties}

The formulation and the solver include $\varepsilon_r$, $K_r$ and $\mu_r$. Two remarks are
needed before using them. First, when $\varepsilon_r\ne1$ a sharp interface without surface
tension is unstable at short wavelengths. The destabilizing normal-field stress grows like
$k$, while gravity does not, so the problem at $Sc=\infty$ is ill-posed at large $k$. A
real miscible interface is regularized at $k\sim1/\delta$, and the smooth-profile problem
is then the relevant one. Second, the effective parameter used in
Ref.~\onlinecite{Traore2024} to predict the effect of the property ratios rescales $T$ with
averaged properties. It cannot account for the jump of the charge density, the jump of
the field or the polarization stress at the interface, which the present formulation
contains.

\subsection{Finite Schmidt number}

For finite $Sc$ the base profile $\bar c$ diffuses, the perturbation obeys the full
Eq.~\eqref{eq:lin_conc}, and the stability problem is non-autonomous. The quasi-steady approximation, with a frozen profile, is accurate only
when the growth rate is large compared with the rate of change of the base state. It fails
close to the threshold \cite{TanHomsy1986,Riaz2006}, which is the regime of interest here.
The natural tool is a non-modal analysis of the propagator over the time-dependent base
state. The present results provide its starting point, since the problem with a frozen
sharp interface is its limit at early times.

\subsection{Realistic stratification}

For a cell of one centimetre and a density contrast of a few percent, $Gr$ is of order
$10^3$ to $10^5$. This is well above the values explored in the simulations, and close to
or above $Gr^*$. In experiments the stationary two-tier mode may therefore be the relevant
onset.

%===========================================================================
\section{Conclusions}
%===========================================================================

We have studied the linear stability of two stacked layers of dielectric liquids with
unipolar injection, in the limit of no molecular diffusion. The kinematic condition at the
interface forbids flow across it in a stationary mode as soon as the interface acts back
on the flow. For equal electrical properties and any stable stratification, the onset of a
single roll spanning both layers is then a Hopf bifurcation. Its threshold grows linearly
with $Gr$ for weak stratification. Above $Gr^*=2137.6$ a stationary mode of two rows of
co-rotating cells, separated by a motionless interface, becomes critical, with
$T_c=477.68$ independent of $Gr$.

The mechanism is not an inertia--gravity oscillation. The interface integrates the
vertical velocity, gravity restores it, and the transit-time lag of the injected charge
sets the frequency and the increase of the threshold. The threshold reported in the
simulations of Ref.~\onlinecite{Traore2024} corresponds to $Gr\approx13$.

The natural continuations are the study of unequal electrical properties with a diffuse
interface, and a non-modal analysis of the time-dependent base state at finite Schmidt
number.

\begin{acknowledgments}
% TODO
\end{acknowledgments}

\section*{Data availability}
The solvers and the data supporting this study are openly available in the repository
\url{https://github.com/pedroavazquez/nonLinearStabilityMonolithic}, in the folder
\texttt{firedrake/two\_layer}.

\appendix

\section{Normal-stress condition}\label{app:normal}

The total stress is $\Sigma=-p\,\mathsf I+\nu(\nabla\bu+\nabla\bu^{T})+M^2(\varepsilon\bE\bE-
\tfrac12\varepsilon E^2\mathsf I)$, whose divergence in each layer gives the viscous,
pressure and Coulomb terms of Eq.~\eqref{eq:ns_layer}. At the displaced interface
$y=\tfrac12+\eta$ the normal stress is continuous. To linear order
\begin{equation}
\jump{\Sigma'_{yy}}+\eta\jump{\bar\Sigma_{yy}'}=0,
\end{equation}
and, by Eq.~\eqref{eq:basestress}, $\jump{\bar\Sigma_{yy}'}=-G$. The perturbation stress is
$\Sigma'_{yy}=-p'+2\nu Dv-M^2\varepsilon\bar E\,D\phi$. The pressure follows from the
horizontal momentum equation, $p'=k^{-2}[-\lambda+\nu(D^2-k^2)]Dv-CM^2\bar q\phi$. With
$\jump{Dv}=0$ the term in $\lambda$ cancels, and Eq.~\eqref{eq:normalstress} follows.

\bibliography{paper}

\end{document}
